\begin{definition}
  For positive $n$, let $\tau(n)$ be the number of positive divisors of $n$.
  The total Lean definition assigns zero at zero.

  \begin{lean}
    def tau (n : ℕ) : ℕ := n.divisors.card
  \end{lean}
\end{definition}

\begin{definition}
  For positive $n$, define $\sigma(n)=\sum_{d\mid n}d$.
  All divisors in this chapter are positive; the divisor set at zero is empty.

  \begin{lean}
    def sigma (n : ℕ) : ℕ := ∑ d ∈ n.divisors, d
  \end{lean}
\end{definition}

\begin{definition}
  For a finitely supported exponent vector $e$, write $V(e)=\prod_p p^{e(p)}$, with the product
  taken over its finite support.

  \begin{lean}
    def factorValue (e : ℕ →₀ ℕ) : ℕ := e.prod (fun p k => p ^ k)
  \end{lean}
\end{definition}

\begin{definition}
  For positive $n$, write $e_n(p)$ for its unique prime multiplicities from Chapter 2.
  The value at zero is assigned the zero vector.

  \begin{lean}
    noncomputable def primeExponents (n : ℕ) : ℕ →₀ ℕ :=
      if hn : 0 < n then (Chapter02.existsUnique_prime_exponents n hn).exists.choose else 0
  \end{lean}
\end{definition}

\begin{theorem}
  \label{ent:thm:dvd-iff-primeexponents-le}
  For positive $d,n$, one has $d\mid n$ precisely when $e_d(p)\le e_n(p)$ for every prime $p$.

  \begin{lean}
    theorem dvd_iff_primeExponents_le (d n : ℕ) (hd : 0 < d) (hn : 0 < n) :
        d ∣ n ↔ primeExponents d ≤ primeExponents n
  \end{lean}
\end{theorem}

\begin{proof}
  If $n=dk$, uniqueness of prime factorization makes $e_n=e_d+e_k$.
  Conversely, subtract the exponent vector of $d$ from that of $n$.
  Its product supplies the quotient witnessing divisibility.

  \begin{lean}
    constructor

    · rintro ⟨k, hk⟩

      have hkpos : 0 < k := by
        nlinarith
      rw [hk, primeExponents_mul d k hd hkpos]
      exact le_add_right le_rfl

    · intro hle

      have hsum : primeExponents d + (primeExponents n - primeExponents d) =
          primeExponents n := add_tsub_cancel_of_le hle

      refine ⟨factorValue (primeExponents n - primeExponents d), ?_⟩

      have hh := congrArg factorValue hsum
      rw [factorValue_add, (primeExponents_spec d hd).2, (primeExponents_spec n hn).2] at hh
      exact hh.symm
  \end{lean}
\end{proof}

\begin{theorem}
  \label{ent:thm:divisor-exponents}
  Every exponent vector $0\le e\le e_n$ gives a positive divisor $V(e)$ of $n$, whose prime
  multiplicities are exactly $e$.
  Consequently these choices describe the divisors uniquely.

  \begin{lean}
    theorem divisor_exponents (n : ℕ) (hn : 0 < n) (e : ℕ →₀ ℕ)
        (he : e ≤ primeExponents n) :
        0 < factorValue e ∧ factorValue e ∣ n ∧ primeExponents (factorValue e) = e
  \end{lean}
\end{theorem}

\begin{proof}
  The support of $e$ lies among the prime factors of $n$, so its product is positive.
  Uniqueness of prime factorization recovers $e$ from that product, and the coordinatewise
  divisibility criterion applies.

  \begin{lean}
    have hp : ∀ p ∈ e.support, Chapter02.isPrime p := by
      intro p hps
      apply (primeExponents_spec n hn).1
      apply Finsupp.mem_support_iff.mpr

      have hpos : 0 < e p := Nat.pos_of_ne_zero (Finsupp.mem_support_iff.mp hps)

      have hle := he p
      omega

    have hpos := factorValue_pos e hp

    have heq := (primeExponents_unique (factorValue e) hpos e hp rfl).symm

    exact ⟨hpos, (dvd_iff_primeExponents_le _ n hpos hn).2 (by
        rw [heq]
        exact he), heq⟩
  \end{lean}
\end{proof}

\begin{theorem}
  \label{ent:thm:tau-eq-prod}
  For positive $n=\prod_p p^{e_n(p)}$, \[ \tau(n)=\prod_{p\mid n}(e_n(p)+1). \] The product is one when $n=1$.

  \begin{lean}
    theorem tau_eq_prod (n : ℕ) (hn : 0 < n) :
        tau n = ∏ p ∈ (primeExponents n).support, (primeExponents n p + 1)
  \end{lean}
\end{theorem}

\begin{proof}
  The divisor-exponent correspondence is a bijection between the divisors and the integer box
  $0\le e(p)\le e_n(p)$.
  Each coordinate has $e_n(p)+1$ choices, so finite counting gives the product.

  \begin{lean}
    classical

    have hcard : n.divisors.card = (Finset.Iic (primeExponents n)).card := by
      apply Finset.card_bij (fun d _ => primeExponents d)

      · intro d hd
        exact Finset.mem_Iic.mpr ((dvd_iff_primeExponents_le d n
          (Nat.pos_of_mem_divisors hd) hn).1 (Nat.dvd_of_mem_divisors hd))

      · intro d hd e he hde

        have hh := congrArg factorValue hde
        rwa [(primeExponents_spec d (Nat.pos_of_mem_divisors hd)).2,
          (primeExponents_spec e (Nat.pos_of_mem_divisors he)).2] at hh

      · intro e he

        obtain ⟨hpos, hdiv, heq⟩ := divisor_exponents n hn e (Finset.mem_Iic.mp he)

        exact ⟨factorValue e, Nat.mem_divisors.mpr ⟨hdiv, ne_of_gt hn⟩, heq⟩
    rw [tau, hcard, Finsupp.card_Iic]
    simp only [Nat.card_Iic]
  \end{lean}
\end{proof}

\begin{definition}
  An integer-valued arithmetic function is \emph{multiplicative} when $f(mn)=f(m)f(n)$ for
  positive coprime $m,n$.
  This convention permits the identically zero function and does not require $f(1)=1$.

  \begin{lean}
    def IsMultiplicative (f : ℕ → ℤ) : Prop :=
      ∀ m n : ℕ, 0 < m → 0 < n → Chapter01.gcd m n = 1 → f (m * n) = f m * f n
  \end{lean}
\end{definition}

\begin{theorem}
  \label{ent:thm:multiplicative-one}
  If a multiplicative function is nonzero at some positive argument, then $f(1)=1$.

  \begin{lean}
    theorem multiplicative_one (f : ℕ → ℤ) (hf : IsMultiplicative f)
        (hne : ∃ n : ℕ, 0 < n ∧ f n ≠ 0) : f 1 = 1
  \end{lean}
\end{theorem}

\begin{proof}
  Choose $n$ with $f(n)\ne0$.
  Multiplicativity at $(1,n)$ gives $f(n)=f(1)f(n)$, and cancellation gives the result.

  \begin{lean}
    obtain ⟨n, hn, hfn⟩ := hne

    have hcop : Chapter01.gcd 1 n = 1 :=
      (Chapter01.gcd_eq_one_iff_exists_mul_add_mul 1 n (Or.inl (by decide))).2
        ⟨1, 0, by ring⟩

    have h := hf 1 n (by decide) hn hcop
    rw [one_mul] at h
    exact mul_right_cancel₀ hfn ((one_mul (f n)).trans h).symm
  \end{lean}
\end{proof}

\begin{theorem}
  \label{ent:thm:existsunique-divisor-split}
  For positive coprime $m,n$, each divisor $d\mid mn$ has a unique expression $d=ab$ with $a\mid
    m$ and $b\mid n$.

  \begin{lean}
    theorem existsUnique_divisor_split (m n d : ℕ) (hm : 0 < m) (hn : 0 < n)
        (hcop : Chapter01.gcd m n = 1) (hd : d ∈ (m * n).divisors) :
        ∃! ab : ℕ × ℕ, ab.1 ∣ m ∧ ab.2 ∣ n ∧ ab.1 * ab.2 = d
  \end{lean}
\end{theorem}

\begin{proof}
  Allocate the prime factors of $d$ to $m$ or $n$ using Euclid\textquotesingle s lemma.
  For uniqueness, if $ab=ce$ with $a,c\mid m$ and $b,e\mid n$, inherited coprimality gives $a\mid
    c$ and $c\mid a$.
  Thus $a=c$, and positivity permits cancellation to obtain $b=e$.
  The same Bezout argument proves that divisors of coprime integers remain coprime.

  \begin{lean}
    obtain ⟨a, b, ha, hb, hab⟩ := exists_divisor_split m n d
      (Nat.pos_of_mem_divisors hd) (Nat.dvd_of_mem_divisors hd)

    refine ⟨(a, b), ⟨ha, hb, hab⟩, ?_⟩
    rintro ⟨c, e⟩ ⟨hc, he, hce⟩

    obtain ⟨hca, heb⟩ := divisor_pair_injective m n hm hn hcop c e a b hc he ha hb
      (hce.trans hab.symm)

    exact Prod.ext hca heb
  \end{lean}
\end{proof}

\begin{definition}
  For an integer-valued arithmetic function $f$, its divisor sum is $F(n)=\sum_{d\mid n}f(d)$.

  \begin{lean}
    def divisorSum (f : ℕ → ℤ) (n : ℕ) : ℤ := ∑ d ∈ n.divisors, f d
  \end{lean}
\end{definition}

\begin{theorem}
  \label{ent:thm:divisorsum-multiplicative}
  If $f$ is multiplicative, then its divisor sum is multiplicative.

  \begin{lean}
    theorem divisorSum_multiplicative (f : ℕ → ℤ) (hf : IsMultiplicative f) :
        IsMultiplicative (divisorSum f)
  \end{lean}
\end{theorem}

\begin{proof}
  For coprime $m,n$, uniquely write every divisor of $mn$ as $ab$ with $a\mid m$ and $b\mid n$.
  Since $a,b$ are coprime, replace $f(ab)$ by $f(a)f(b)$ and factor the double sum.

  \begin{lean}
    intro m n hm hn hcop
    rw [sum_divisors_mul f m n hm hn hcop, divisorSum, divisorSum, Finset.sum_mul_sum]
    apply Finset.sum_congr rfl
    intro a ha
    apply Finset.sum_congr rfl
    intro b hb
    exact hf a b (Nat.pos_of_mem_divisors ha) (Nat.pos_of_mem_divisors hb)
      (gcd_divisors_eq_one m n a b hm hn hcop (Nat.dvd_of_mem_divisors ha)
        (Nat.dvd_of_mem_divisors hb))
  \end{lean}
\end{proof}

\begin{theorem}
  \label{ent:thm:tau-multiplicative}
  The function $\tau$ is multiplicative.

  \begin{lean}
    theorem tau_multiplicative : IsMultiplicative (fun n => (tau n : ℤ))
  \end{lean}
\end{theorem}

\begin{proof}
  Take the divisor sum of the constant function one, which is multiplicative.

  \begin{lean}
    have hh := divisorSum_multiplicative (fun _ => 1) (by
      intro m n hm hn hcop
      ring)
    simpa only [IsMultiplicative, divisorSum, Finset.sum_const, nsmul_eq_mul, mul_one, tau] using hh
  \end{lean}
\end{proof}

\begin{theorem}
  \label{ent:thm:sigma-multiplicative}
  The function $\sigma$ is multiplicative.

  \begin{lean}
    theorem sigma_multiplicative : IsMultiplicative (fun n => (sigma n : ℤ))
  \end{lean}
\end{theorem}

\begin{proof}
  Take the divisor sum of $f(d)=d$, which is multiplicative.

  \begin{lean}
    have hh := divisorSum_multiplicative (fun n => (n : ℤ)) (by
      intro m n hm hn hcop
      exact Nat.cast_mul m n)
    simpa only [IsMultiplicative, divisorSum, sigma, Nat.cast_sum] using hh
  \end{lean}
\end{proof}

\begin{theorem}
  \label{ent:thm:multiplicative-prime-powers}
  For $n>1$, a multiplicative function satisfies $f(n)=\prod_{p\mid n}f(p^{e_n(p)})$.
  This also holds for the zero function.

  \begin{lean}
    theorem multiplicative_prime_powers (f : ℕ → ℤ) (hf : IsMultiplicative f)
        (n : ℕ) (hn : 1 < n) :
        f n = ∏ p ∈ (primeExponents n).support, f (p ^ primeExponents n p)
  \end{lean}
\end{theorem}

\begin{proof}
  Distinct prime powers are coprime.
  Apply multiplicativity successively to the nonempty product of prime powers.
  Nonemptiness avoids making an unwarranted assumption that $f(1)=1$.

  \begin{lean}
    have hnpos : 0 < n := by
      omega

    have hspec := primeExponents_spec n hnpos

    have hs : (primeExponents n).support.Nonempty := by
      by_contra h

      have he := Finset.not_nonempty_iff_eq_empty.mp h

      have hv := hspec.2
      change (∏ p ∈ (primeExponents n).support, p ^ primeExponents n p) = n at hv
      rw [he, Finset.prod_empty] at hv
      omega
    calc
      f n = f (factorValue (primeExponents n)) := congrArg f hspec.2.symm
      _ = _ := multiplicative_prod _ hs (fun p => p ^ primeExponents n p)
        (fun p hp => pow_pos (by
            have := (hspec.1 p hp).1
            omega) _)
        (fun p hp q hq hne => gcd_prime_powers p q _ _ (hspec.1 p hp) (hspec.1 q hq) hne) f hf
  \end{lean}
\end{proof}

\begin{theorem}
  \label{ent:thm:tau-prime-pow}
  If $p$ is prime and $k\ge0$, then $\tau(p^k)=k+1$.

  \begin{lean}
    theorem tau_prime_pow (p k : ℕ) (hp : Chapter02.isPrime p) : tau (p ^ k) = k + 1
  \end{lean}
\end{theorem}

\begin{proof}
  The divisors are exactly $1,p,\ldots,p^k$, and these powers are distinct.

  \begin{lean}
    classical
    rw [tau, divisors_prime_pow p k hp, Finset.card_image_of_injective,
      Finset.card_range]

    exact pow_right_injective₀ (by have := hp.1; omega : 0 < p) (by have := hp.1; omega : p ≠ 1)
  \end{lean}
\end{proof}

\begin{theorem}
  \label{ent:thm:sigma-prime-pow}
  If $p$ is prime and $k\ge0$, then $\sigma(p^k)=1+p+\cdots+p^k$.

  \begin{lean}
    theorem sigma_prime_pow (p k : ℕ) (hp : Chapter02.isPrime p) :
        sigma (p ^ k) = ∑ j ∈ Finset.range (k + 1), p ^ j
  \end{lean}
\end{theorem}

\begin{proof}
  Sum the elements of the prime-power divisor list.

  \begin{lean}
    classical
    rw [sigma, divisors_prime_pow p k hp, Finset.sum_image]
    intro i hi j hj h
    exact pow_right_injective₀ (by have := hp.1; omega : 0 < p) (by have := hp.1; omega : p ≠ 1) h
  \end{lean}
\end{proof}

\begin{theorem}
  \label{ent:thm:sigma-eq-prod}
  For $n>1$, \[ \sigma(n)=\prod_{p\mid n}\bigl(1+p+\cdots+p^{e_n(p)}\bigr). \]

  \begin{lean}
    theorem sigma_eq_prod (n : ℕ) (hn : 1 < n) :
        (sigma n : ℤ) = ∏ p ∈ (primeExponents n).support,
          ((∑ j ∈ Finset.range (primeExponents n p + 1), p ^ j : ℕ) : ℤ)
  \end{lean}
\end{theorem}

\begin{proof}
  Apply the prime-power formula for multiplicative functions to $\sigma$, then substitute the
  geometric sum for each prime-power value.

  \begin{lean}
    rw [multiplicative_prime_powers _ sigma_multiplicative n hn]
    apply Finset.prod_congr rfl
    intro p hp
    rw [sigma_prime_pow p _ ((primeExponents_spec n (by omega)).1 p hp)]
  \end{lean}
\end{proof}

\begin{theorem}
  \label{ent:thm:sigma-eq-prod-quotients}
  For $n>1$, \[ \sigma(n)=\prod_{p\mid n}\frac{p^{e_n(p)+1}-1}{p-1}. \] Each quotient is an integer.

  \begin{lean}
    theorem sigma_eq_prod_quotients (n : ℕ) (hn : 1 < n) :
        (sigma n : ℤ) = ∏ p ∈ (primeExponents n).support,
          (((p ^ (primeExponents n p + 1) - 1) / (p - 1) : ℕ) : ℤ)
  \end{lean}
\end{theorem}

\begin{proof}
  The geometric identity $(p-1)(1+p+\cdots+p^k)=p^{k+1}-1$ follows by induction.
  Replace each geometric sum in the product by this exact quotient.

  \begin{lean}
    rw [sigma_eq_prod n hn]
    apply Finset.prod_congr rfl
    intro p hp

    rw [geometric_sum_eq_div p _ ((primeExponents_spec n (by omega)).1 p hp).1]
  \end{lean}
\end{proof}

\begin{theorem}
  \label{ent:thm:tau-sigma-one}
  One has $\tau(1)=\sigma(1)=1$.

  \begin{lean}
    theorem tau_sigma_one : tau 1 = 1 ∧ sigma 1 = 1
  \end{lean}
\end{theorem}

\begin{proof}
  The only positive divisor of one is one.

  \begin{lean}
    simp only [tau, sigma, Nat.divisors_one, Finset.card_singleton, Finset.sum_singleton, and_self]
  \end{lean}
\end{proof}

\begin{theorem}
  \label{ent:thm:tau-eq-two-iff}
  For $n>1$, one has $\tau(n)=2$ exactly when $n$ is prime.

  \begin{lean}
    theorem tau_eq_two_iff (n : ℕ) (hn : 1 < n) : tau n = 2 ↔ Chapter02.isPrime n
  \end{lean}
\end{theorem}

\begin{proof}
  The divisors already include the distinct elements one and $n$.
  If there are only two, these are all the divisors.
  Conversely, the definition of primality says precisely that no other divisor exists.

  \begin{lean}
    have hnpos : 0 < n := by
      omega

    have hnne : n ≠ 0 := ne_of_gt hnpos

    have hsub : ({1, n} : Finset ℕ) ⊆ n.divisors := by
      intro d hd
      simp only [Finset.mem_insert, Finset.mem_singleton] at hd
      rcases hd with he | he

      · rw [he]
        exact Nat.one_mem_divisors.mpr (ne_of_gt hnpos)

      · rw [he]
        exact Nat.mem_divisors_self n (ne_of_gt hnpos)

    have hcard : ({1, n} : Finset ℕ).card = 2 := by
      rw [Finset.card_insert_of_notMem, Finset.card_singleton]
      simpa only [Finset.mem_singleton] using (by omega : (1 : ℕ) ≠ n)

    constructor

    · intro ht

      have heq : n.divisors = {1, n} :=
        (Finset.eq_of_subset_of_card_le hsub (by
            change tau n ≤ _
            rw [ht, hcard])).symm

      refine ⟨hn, ?_⟩
      intro d hd

      have hmem := Nat.mem_divisors.mpr ⟨hd, ne_of_gt hnpos⟩
      rw [heq, Finset.mem_insert, Finset.mem_singleton] at hmem
      exact hmem

    · intro hp

      have heq : n.divisors = {1, n} := by
        apply Finset.Subset.antisymm _ hsub
        intro d hd
        exact Finset.mem_insert.mpr ((hp.2 d (Nat.dvd_of_mem_divisors hd)).imp_right
          Finset.mem_singleton.mpr)

      rw [tau, heq, hcard]
  \end{lean}
\end{proof}

\begin{theorem}
  \label{ent:thm:sigma-eq-succ-iff}
  For $n>1$, one has $\sigma(n)=n+1$ exactly when $n$ is prime.

  \begin{lean}
    theorem sigma_eq_succ_iff (n : ℕ) (hn : 1 < n) : sigma n = n + 1 ↔ Chapter02.isPrime n
  \end{lean}
\end{theorem}

\begin{proof}
  One and $n$ already contribute $n+1$.
  Any further positive divisor would strictly increase the sum.
  Conversely, the two divisors of a prime sum to $n+1$.

  \begin{lean}
    have hnpos : 0 < n := by
      omega

    have hnne : n ≠ 0 := ne_of_gt hnpos

    constructor

    · intro hs
      refine ⟨hn, ?_⟩
      intro d hd
      by_cases hd1 : d = 1

      · exact Or.inl hd1

      · by_cases hdn : d = n

        · exact Or.inr hdn

        · have hmem := Nat.mem_divisors.mpr ⟨hd, ne_of_gt hnpos⟩

          have hdpos := Nat.pos_of_mem_divisors hmem

          have hsub : ({d, 1, n} : Finset ℕ) ⊆ n.divisors := by
            intro k hk
            simp only [Finset.mem_insert, Finset.mem_singleton] at hk
            rcases hk with he | he | he

            · rwa [he]

            · rw [he]
              exact Nat.one_mem_divisors.mpr (ne_of_gt hnpos)

            · rw [he]
              exact Nat.mem_divisors_self n (ne_of_gt hnpos)

          have hle := Finset.sum_le_sum_of_subset hsub (f := fun k => k)
          rw [Finset.sum_insert (by simp only [Finset.mem_insert, Finset.mem_singleton,
            hd1, hdn, or_self, not_false_eq_true]),
            Finset.sum_insert (by simpa only [Finset.mem_singleton] using
              (by omega : (1 : ℕ) ≠ n)), Finset.sum_singleton] at hle
          change d + (1 + n) ≤ sigma n at hle
          omega

    · intro hp

      have heq : n.divisors = {1, n} := by
        ext d
        constructor

        · intro hd
          exact Finset.mem_insert.mpr ((hp.2 d (Nat.dvd_of_mem_divisors hd)).imp_right
            Finset.mem_singleton.mpr)

        · intro hd
          rcases Finset.mem_insert.mp hd with he | he

          · rw [he]
            exact Nat.one_mem_divisors.mpr hnne

          · rw [Finset.mem_singleton.mp he]
            exact Nat.mem_divisors_self n hnne

      rw [sigma, heq, Finset.sum_insert (by simpa only [Finset.mem_singleton] using
        (by omega : (1 : ℕ) ≠ n)), Finset.sum_singleton]
      omega
  \end{lean}
\end{proof}

\begin{theorem}
  \label{ent:thm:prod-divisors-sq}
  For positive $n$, the product of its divisors satisfies \[ \left(\prod_{d\mid n}d\right)^2=n^{\tau(n)}. \] Equivalently, as positive real numbers, the product is $n^{\tau(n)/2}$; the squared form avoids truncating the half-exponent.

  \begin{lean}
    theorem prod_divisors_sq (n : ℕ) (hn : 0 < n) :
        (∏ d ∈ n.divisors, d) ^ 2 = n ^ tau n
  \end{lean}
\end{theorem}

\begin{proof}
  The map $d\mapsto n/d$ permutes the positive divisors.
  Multiply the original divisor product by this permuted product; each paired factor is $d(n/d)=n$.

  \begin{lean}
    have hperm : (∏ d ∈ n.divisors, n / d) = ∏ d ∈ n.divisors, d := by
      apply Finset.prod_bij (fun d _ => n / d)

      · intro d hd
        exact Nat.mem_divisors.mpr
          ⟨Nat.div_dvd_of_dvd (Nat.dvd_of_mem_divisors hd), ne_of_gt hn⟩

      · intro d hd e he hde

        have h := congrArg (fun k => n / k) hde
        rwa [Nat.div_div_self (Nat.dvd_of_mem_divisors hd) (ne_of_gt hn),
          Nat.div_div_self (Nat.dvd_of_mem_divisors he) (ne_of_gt hn)] at h

      · intro d hd
        refine ⟨n / d, Nat.mem_divisors.mpr
          ⟨Nat.div_dvd_of_dvd (Nat.dvd_of_mem_divisors hd), ne_of_gt hn⟩, ?_⟩

        exact Nat.div_div_self (Nat.dvd_of_mem_divisors hd) (ne_of_gt hn)

      · intro d hd
        rfl

    rw [pow_two]
    nth_rw 2 [← hperm]
    rw [← Finset.prod_mul_distrib]

    have hterm : ∀ d ∈ n.divisors, d * (n / d) = n := by
      intro d hd
      exact Nat.mul_div_cancel' (Nat.dvd_of_mem_divisors hd)

    rw [Finset.prod_congr rfl hterm, Finset.prod_const]
    rfl
  \end{lean}
\end{proof}
